% Marker for text the team must still replace before submission. Search the source for \needchange.
\providecommand{\needchange}[1]{\textcolor{red}{\textbf{[NEED TO CHANGE: #1]}}}

\section{Three connected research questions and verified gap}

\noindent
Organizations increasingly face scarce GPU capacity as AI workloads grow. First-come-first-served allocation ignores project value and urgency, potentially misallocating scarce compute and leaving carbon costs unaccounted for.

\vspace{3pt}

\noindent
The closest work is Zachariadis et al.'s FAIR-Compute report~\cite{fair_compute2026}. Its method is a trace-driven simulation: it replays a public academic workload (Purdue Anvil, about 20.9 million jobs) under FIFO, Slurm-style priority, a tuned multi-objective priority, and an offline optimization benchmark. Its evidence is that online, value-blind rules recover only about 65\% of the available scientific value (1552--1612 of 2388 units; their Section~5.4), so roughly a third is lost. It also marks the two gaps we address: it does not quantify carbon (Section~3.6), and it warns that self-declared value will be exaggerated (Section~4.2) while leaving mechanisms under which truthful reporting is optimal to future work (Section~7). We build that missing piece. We trade their trace-level scale for a six-team instance that is solved exactly, so the allocation, the payments, and the truthfulness claim can each be re-checked by a reader.

\vspace{3pt}

\noindent
\textbf{Q1---Economics:} How should an institution allocate scarce GPU-hours to maximize project value net of carbon cost, when teams privately know their own value and can misreport it?\par
\textbf{Q2---Computation:} Can an exactly-solvable multi-team VCG mechanism allocate more efficiently than first-come-first-served, while staying auditable?\par
\textbf{Q3---Behavior:} Do participants report values close to the truthful benchmark, or does behavior deviate---and do their predictions and reflections show understanding of the rule?

\vspace{3pt}

\noindent
Together: can a transparent, carbon-aware rule replace arrival order, and does its advantage survive behavioral and computational limits? Figure~\ref{fig:teaser} states the three questions in its top row and traces the model (Q1) and its exact computation (Q2) below them; the behavioral test (Q3) is described in Section 4.

\section{Answer to Q1: incentives and strategic model}

\noindent
\textbf{Hypothesis.} An incentive-compatible, carbon-aware mechanism increases realized value and reduces emissions per unit of value, without misreporting.

\vspace{3pt}

\noindent
\textbf{Model.} Six teams $i\in\{A,\dots,F\}$ compete for $K=100$ GPU-hours. Each privately knows value $v_i \in \{3,6,9\}$; demand $d_i$ is observable, emissions $e_i=0.24d_i$ kg CO$_2$e, carbon weight $\lambda=0.5$ (stated parameters, not empirical estimates). Teams simultaneously submit one report $r_i$ (static, incomplete information).

\vspace{3pt}

\noindent
\textbf{Allocation and payment.} Score $s_i=r_i-\lambda e_i$; select feasible $S^*$ maximizing $\sum_{i\in S}s_i$ s.t. $\sum_{i\in S}d_i\le K$ (demand enters only as the constraint). Each winner pays the Clarke--Groves externality $p_i = W_{-i}^* - \sum_{j\in S^*\setminus\{i\}}s_j$~\cite{clarke1971,groves1973} (Appendix~A); utility is $u_i=v_i-\lambda e_i-p_i$ if selected, else 0.

\vspace{3pt}

\noindent
\textbf{Intuition.} The payment $p_i$ is the score the other teams lose because $i$ is selected, and it does not depend on $r_i$. Substituting it gives $u_i = (\text{true total score of } S^*) - W_{-i}^*$; the second term is fixed, so a team can raise its utility only by letting the rule pick the truly best group, which is what a truthful report does. A report therefore matters only through whether it crosses the team's selection boundary. Team C ($v_C=3$) is selected once $r_C\ge 4.2$, and every report above that boundary pays the same 2.4 and yields $u_C=-1.2$: overreporting more never costs more, it only buys a seat that is worth less than its price (Appendix~F.1).

\vspace{3pt}

\noindent
\textbf{Solution concept.} Truthful reporting $r_i=v_i$ is dominant---stronger than Nash/Bayesian Nash equilibrium~\cite{nash1950,harsanyi1968}---holding only because each team bears its own carbon penalty in $u_i$ (else the optimum shifts to $r_i=v_i+\lambda e_i$) and credits carry real future cost; verified numerically over $[0,10]$ for every team and for every carbon weight $\lambda\in\{0,0.1,\dots,1\}$ (Appendix~A.2). A reduced $2\times2$ payoff matrix for Teams C and E is in Appendix~A.1 (Table~\ref{tab:matrix}).

\vspace{3pt}

\noindent
\textbf{Three lenses.} Game theory yields the dominant-strategy prediction. Social choice frames the objective as value net of carbon under capacity and surfaces two trade-offs: efficiency against access (VCG serves 4 teams against an FCFS mean of 4.93), and truthfulness against budget balance (the 96 credits go to the organization rather than back to teams; returning them would in general break truthfulness~\cite{greenlaffont1977}, so they are the price of a dominant strategy). Mechanism design specifies the report space, the rules, and the verified tests (Appendix~C).

\section{Answer to Q2: computation, auction comparison, and artifacts}

\noindent
\textbf{Inputs.} Six teams (A--F): demands $(25,20,15,25,20,15)$, values $(9,6,3,9,6,3)$. FCFS is evaluated over all 720 arrival orders; VCG enumerates all 64 feasible subsets and re-solves per winner for payments (Appendix~A.1--A.2).

\textbf{Result.} VCG selects $\{A,B,D,E\}$ (score 19.2)---no FCFS order of 720 exceeds it (range 15.6–19.2; mean 16.96). Emissions per value fall from 0.81--0.84 to 0.72 kg CO$_2$e. One verified synthetic instance, not general evidence (Appendix~A, GitHub).

\vspace{3pt}

\noindent
\textbf{Parameter change.} Sweeping $\lambda$ from 0 to 1 in steps of 0.1 changes payoffs but not strategy. The selected group is $\{A,B,D,E\}$ for every $\lambda\in[0.1,1]$. The payment per winner falls from 5.28 ($\lambda=0.1$) to 2.4 ($\lambda=0.5$) and to 0 (from $\lambda=0.9$), because the low-value teams stop competing for a seat; the utility of A and D is 3.12, 3.6, and 3.0 at $\lambda=0.1, 0.5, 1$. Truthful reporting is optimal at all eleven values. $\lambda$ is therefore a policy weight on carbon, not an incentive parameter, and it cannot be chosen by maximizing project value, since on value alone $\lambda=0$ always wins. A second change splits capacity into two GPU types (H100, 50 h at 0.24 kg; A100, 50 h at 0.14 kg): the same four teams win, the two largest jobs move to the cleaner pool, emissions fall from 21.6 to 16.6 kg, and truthful reporting stays optimal (Appendix~A.2).

\vspace{3pt}

\noindent
\textbf{Analogy boundary.} The auction analogy holds only in part. Payments are non-transferable priority credits rather than money, so the dominant-strategy result needs credits to carry a real future opportunity cost, which the code cannot verify. True values are never observed: allocation and payment use reports only. And unlike a single-item auction, several teams win under a capacity (knapsack) constraint that is solved by exact search over 64 groups; only with one slot and $\lambda=0$ does the rule reduce to a Vickrey second-price auction~\cite{vickrey1961}, a case our code reproduces (the winner pays the second-highest value).

\vspace{3pt}

\noindent
\textbf{Artifacts.} GitHub: \GitHubLink{} (commit \texttt{8a656a4}). Notebook: \NotebookLink. Hugging Face: \HuggingFaceLink{} (Space commit \texttt{ad5604b}). A0 poster: \texttt{PS2-FP05-TeamA-A0-Poster.pdf}.

\section{Answer to Q3: behavior and interdisciplinary integration}

\noindent
\textbf{Rational benchmark.} Section~2 establishes $r_i=v_i$ as dominant under VCG; FCFS has no reporting decision at all.

\vspace{3pt}

\noindent
\textbf{Behavioral mechanism and competing explanation.} Laboratory second-price auctions show persistent overbidding even though truthful bidding is dominant~\cite{kagel1987}. Two explanations predict the same overreport in our setting. (i)~\emph{Comprehension:} a sealed-bid VCG rule is not obviously strategy-proof, so players may not see that the truth is optimal~\cite{li2017}. (ii)~\emph{Preference:} players understand the rule but value winning itself~\cite{cooperfang2008}.

\vspace{3pt}

\noindent
\textbf{Computed consequences of each behavior.} The model shows what each deviation would cost; these rows are computed, not observed (rivals truthful, $\lambda=0.5$).

{\footnotesize
\begin{tabularx}{\columnwidth}{@{}p{2.25cm}p{1.45cm}Y@{}}\toprule
\textbf{Behavior} & \textbf{Own utility} & \textbf{Effect on other teams}\\\midrule
Truthful (all) & A,D 3.6; B,E 1.2; C,F 0 & baseline\\
A overreports, $9\to10$ & 3.6 (same) & none: A already holds a seat and its price is unchanged\\
C overreports, $3\to4.5$ & $0\to-1.2$ & E loses its seat; B falls to 0; A, D fall to 2.4\\
B underreports, $6\to4$ & $1.2\to0$ & B loses its seat; C and F enter\\\bottomrule
\end{tabularx}}

\vspace{3pt}

\noindent
A deviation in either direction never helps the deviating team, but an overreport that crosses the boundary also harms others: the rule guarantees that the liar does not gain, not that bystanders are protected.

\vspace{3pt}

\noindent
\textbf{Observed evidence (exploratory).} We recorded $n=5$ completed plays of the Space (\texttt{hf-space/plays.csv}), played by the three authors on October~3, 2026. Each player held a private value, chose a request and a report, and predicted the outcome before seeing it; the five rivals were simulated and truthful, and stakes were hypothetical.

{\footnotesize
\begin{tabularx}{\columnwidth}{@{}Ycc@{}}\toprule
\textbf{Measure} & \textbf{Benchmark} & \textbf{Observed}\\\midrule
Truthful reports ($r_i=v_i$) & 5/5 & 2/5\\
Mean $|r_i-v_i|$ & 0 & 1.8\\
Overreports / underreports & 0 / 0 & 3 / 0\\
Selection predicted correctly & --- & 3/5\\
Overreport with a wrong prediction & 0 & 1\\
Overreport with a correct prediction & 0 & 2\\
Plays in which a misreport raised utility & 0 & 0\\\bottomrule
\end{tabularx}}

\vspace{3pt}

\noindent
\emph{Agreement.} The payoff prediction held in every play: no overreport raised the player's utility. Two overreports left utility unchanged, and one bought a seat that cost 36 credits and was worth nothing net of carbon ($u=-3.6$ against 0 for the truth). \emph{Gap.} The behavioral prediction did not hold: three of five players overreported (by 4, 4, and 1), although the report field starts at the true value. Both wrong predictions were pessimistic (``left out'' when the team was selected). \emph{Consequence.} We retain the dominant-strategy result as a statement about payoffs and qualify the assumption that players see it: truthful reporting was optimal but not typical. The written reflections were left blank in all five plays. Asked afterwards, the players gave two reasons for overreporting: wanting to be selected, and misunderstanding the rule. These match explanations~(ii) and~(i), but they are recollections given after the result was known and were not recorded per play, so we do not assign them to individual plays. One further observation is consistent with explanation~(i): a symposium reviewer who had read the rule asked whether truthful reporting is optimal with and without the carbon charge (Appendix~D). These are five self-plays by the authors, who know the rule, against simulated rivals; they are exploratory and support no population claim. These observations motivate a larger classroom sample with written reflections, but they do not predict how often other participants will misreport.

\vspace{3pt}

\noindent
\textbf{Interaction.} Users play one team through both FCFS and VCG stages, predicting their own selection before seeing the result, then inspect the ranking, payments, and utility (full mechanics in Appendix~F.1). All evidence is exploratory, with participant count and conditions reported.

\vspace{3pt}

\noindent
\textbf{Discriminating observation.} The prediction step separates the two explanations. An overreport with a \emph{wrong} selection prediction points to comprehension~(i) and would lead us to revise the interface (show the what-if curve before the report). An overreport with a \emph{correct} prediction and a reflection about wanting to win points to preference~(ii) and would lead us to revise the model (add a value of winning to $u_i$). Neither supports a causal claim. The first five plays contain one case of the first kind and two overreports with a correct prediction but no written reflection, so the next step is the interface revision, which also requires the reflection before the result line can be copied (planned).

\section{Advanced development, real-world impact, and roadmap}

\noindent
This proposal builds on two Nobel-recognized foundations: Vickrey's theory of incentives under asymmetric information~\cite{vickrey1961,nobel1996} and Hurwicz--Maskin--Myerson's mechanism design theory~\cite{nobel2007}---agents can be motivated to reveal the truth without verifying private information. The model is specified enough to support a transparent pilot design, but its parameters remain classroom assumptions. Integrating incentive-compatible allocation, an explicit carbon objective, and exploratory behavioral evidence addresses the two gaps left by Section~1's closest prior work.

\vspace{3pt}

\noindent
The concrete case is a shared university GPU cluster. Teams benefit from allocation reflecting true value; administrators gain an auditable record; the public bears the carbon cost regardless of who benefits. Honest reporting is incentivized, but truthfulness relative to a team's actual beliefs is not independently verifiable. The carbon weight $\lambda$ is the institution's choice, for example an internal carbon price converted into value units; our sweep shows the allocation is stable for $\lambda\in[0.1,1]$, while a weight that is too heavy stops research (only A and D are served above $\lambda=1.25$, and no team above $\lambda=1.5$). Academic adoption would accept a transparent pilot; industry would demand robustness at hundreds of teams; public institutions would require an independently re-runnable audit. Next test: a pilot comparing self-reported value against ex-post outcomes at a scale where exact solving is no longer tractable. Planned voting bridge: a variant in which teams rank the projects and the rule aggregates the rankings (plurality against Borda), to compare manipulation under voting with misreporting under VCG; this is planned work, not a result.

\vspace{3pt}

\begin{table}[h]
\caption{From foundations to a collaborative interdisciplinary contribution.}
\begin{tabularx}{\columnwidth}{@{}p{1.6cm}Y@{}}
\toprule
\textbf{Horizon}&\textbf{Milestone and evidence}\\
\midrule
Foundation&Vickrey~\cite{nobel1996} and Hurwicz--Maskin--Myerson~\cite{nobel2007}: incentive-compatible allocation without verifying private information.\\
PS2&Six-team carbon-aware VCG mechanism, exactly solved and verified against a random-order first-come-first-served baseline over all 720 arrival orders (Appendix~A).\\
Symposium&Poster on institutional GPU governance (Sept.~28). Peer challenges: carbon-weight calibration, a single GPU type, missing play data, and whether truth stays optimal with a carbon charge. Response: a $\lambda$ sweep, a two-GPU-pool extension, computed behavior scenarios, and five recorded self-plays (Appendix~D).\\
Next study&A rerun of the FCFS--VCG comparison over randomly drawn values and requests, then a pilot comparing self-reported value against independent, ex-post project outcomes at a scale where exact solving is no longer tractable.\\
2056&Auditable mechanism-design tool adopted by a real institution for actual, not simulated, compute governance.\\
\bottomrule
\end{tabularx}
\end{table}

\vspace{3pt}

\noindent
\noindent\textbf{Expected 2056 Nobel Prize and Turing Award contribution statement.}
By 2056, we aim to be recognized for an efficient, publicly auditable governance framework for scarce shared infrastructure, integrating an incentive-compatible mechanism, verifiable computation, and behavioral evidence on truthful reporting. Progress moves from this six-team simulation to a real institutional pilot, correctable through a public, re-runnable audit log.

\subsection*{Open Science Statement}
GitHub (\url{https://github.com/dku-comsci-econ206-Autumn2026/FP05-PS2-Team-A}) contains the FCFS/VCG code, tests, and outputs, including the carbon-weight sweep (\texttt{lambda\_sweep.csv}); a self-contained Colab notebook (\url{https://colab.research.google.com/drive/14eo8a2HWUbbBtv-m-ia4xRiL4V6Iwtun}) that also holds the two-GPU-pool extension; a Hugging Face Space (\url{https://huggingface.co/spaces/dku-comsci-econ206-2026/PS2_Team_A}, commit \texttt{ad5604b}); README with run instructions. Inputs are synthetic ($d_i\in\{15,20,25\}$, $v_i\in\{3,6,9\}$; the Space allows wider ranges), clearly labeled. Review-ready results come from commit \texttt{d93fc2e}; final results from commit \texttt{8a656a4}. Limitation: exact enumeration does not scale past small team counts.

\subsection*{Statement of Contribution to the UN's SDGs}
Connects to SDG~13 (carbon priced into allocation) and SDG~9 (auditable infrastructure rule). Audience: compute administrators and research teams; minimal access needs. Indicator: emissions per unit of allocated value. Possible unintended effect: urgency inflation beyond the report's resolution, or $\lambda$ disadvantaging compute-heavy research. Intended mechanism only---no adoption has occurred.
